
\section{Open Research Directions}
\label{sec:discussion:threads}

\subsection{Closing the Modeling and Deployment Gap}
\label{sec:discussion:modeling}

Every system in this thesis was developed and primarily evaluated in simulation, with real-world experiments serving as validation rather than the primary test environment.
The sim-to-real gap is a pervasive problem: physics simulators approximate contact dynamics imperfectly, and discrepancies accumulate over long planning horizons.
\textit{PokeRRT} addressed this pragmatically through replanning---treating the gap as disturbance and replanning when observations diverge from predictions.
But replanning is a band-aid; closing the gap requires either better simulation fidelity or learned residual models that correct simulator errors from real interaction data.
Future work should explore differentiable physics engines that allow gradient-based refinement from real-world trajectory data, and online adaptation mechanisms that update dynamics models during deployment.
The generative methods raise a related challenge: their training distributions were generated by specific planners in specific environments.
Understanding how coverage and composition of training data affects generalization---and developing systematic approaches to training data design for dynamics-aware generative models---is an open problem with direct implications for deployment reliability.

\subsection{Scaling with Data and New Sensing Modalities}
\label{sec:discussion:datascaling}

The generative methods in this thesis were trained on datasets that are large by robotics standards but small by machine learning standards.
Whether the quality of dynamics-aware trajectory generation scales with data volume---in the manner of language model scaling laws---is unknown.
If it does, the practical bottleneck becomes data collection; if it does not, the architecture or training objective requires rethinking.
Separately, the representations used for dynamics conditioning in this thesis are relatively simple: scalar payload mass, binary embodiment vectors, discrete constraint-satisfying samples.
New sensing modalities---high-resolution tactile skins, whole-body force-torque sensing, high-speed event cameras---could provide richer constraint representations that enable conditioning on more complex and subtle dynamics.
Understanding which sensing modalities are most informative for dynamics-constrained generation, and how to incorporate them efficiently into conditioning mechanisms, is a promising direction.

\subsection{Formal Guarantees and Safety in Learned Dynamics Systems}
\label{sec:discussion:guarantees}

None of the generative methods in this thesis come with formal guarantees.
Payload diffusion achieves 67.6\% workspace coverage at $3\times$ nominal payload, but the inaccessible portion of workspace is empirically characterized, not formally bounded.
\textit{DEFT} maintains high success rates under out-of-distribution failures, but the conditions under which it fails are not analytically characterizable.
For deployment in safety-critical applications---surgical assistance, unstructured field environments, manipulation near humans---probabilistic success rates are insufficient.
The fundamental challenge is that the guarantees classical control theory provides (Lyapunov stability, energy-based safety certificates, control barrier functions) are derived analytically from known dynamics, while learned models are defined implicitly through their weights.
Promising directions include training diffusion models with energy-based objectives that respect known conservative properties of the dynamics, verifying constraint satisfaction post-hoc through lightweight formal checkers applied to generated trajectories, and hybrid architectures that use learned models for initialization and certified optimization for final execution.

\subsection{Hardware Co-Design and Cross-Embodiment Generalization}
\label{sec:discussion:actuation}

This thesis assumed standard articulated manipulators.
But several contributions push against the limits of what standard hardware can do: payload diffusion operates at the edge of rated joint torques, poking involves high-speed end-effector impacts that cause hardware wear, and whole-body manipulation requires arm surfaces not normally treated as manipulation surfaces.
A natural question is whether the dynamic capabilities explored here call for specialized actuation---compliant joints, variable-stiffness actuators, high-speed lightweight wrists---or whether general-purpose hardware, paired with better motion generation, suffices.
This is a co-design question that the robotics community has not systematically addressed in the context of dynamics-constrained planning.
A related question concerns cross-embodiment generalization.
\textit{DEFT} treats each failure configuration as a distinct embodiment; training one model to generalize across robot platforms, arm configurations, and failure modes remains an open challenge.
Addressing it would substantially reduce the per-platform engineering cost of deploying dynamics-aware systems.